\documentclass{standalone}
\usepackage{tikz}
\usetikzlibrary{shapes.geometric, arrows}

\begin{document}
\begin{tikzpicture}[node distance=2cm]
    % Definir nodos del flujo
    \node (start) [draw, ellipse, fill=green!20] {Inicio};
    \node (cities) [draw, rectangle, fill=blue!20, below=of start] {Generación de Ciudades};
    \node (population) [draw, rectangle, fill=blue!20, below=of cities] {Creación de Población Inicial};
    \node (fitness) [draw, rectangle, fill=yellow!20, below=of population] {Evaluación de Fitness};
    \node (selection) [draw, rectangle, fill=orange!20, below=of fitness] {Selección};
    \node (crossover) [draw, rectangle, fill=orange!20, below=of selection] {Cruza};
    \node (mutation) [draw, rectangle, fill=red!20, below=of crossover] {Mutación};
    \node (newpop) [draw, rectangle, fill=red!20, below=of mutation] {Generación de Nueva Población};
    \node (check) [draw, diamond, fill=purple!20, below=of newpop] {¿Máx Generaciones?};
    \node (visual) [draw, rectangle, fill=green!20, right=of check] {Visualización de Resultados};
    \node (end) [draw, ellipse, fill=red!20, below=of check] {Fin};

    % Dibujar flechas del flujo
    \draw[->] (start) -- (cities);
    \draw[->] (cities) -- (population);
    \draw[->] (population) -- (fitness);
    \draw[->] (fitness) -- (selection);
    \draw[->] (selection) -- (crossover);
    \draw[->] (crossover) -- (mutation);
    \draw[->] (mutation) -- (newpop);
    \draw[->] (newpop) -- (check);
    \draw[->] (check) -- node[above] {No} (fitness);
    \draw[->] (check) -- node[right] {Sí} (visual);
    \draw[->] (visual) -- (end);
\end{tikzpicture}
\end{document}